\scanpage{303}
\mspage{\textcircled{2}}
Un objet de $A_*$ sera généralement noté par un symbole de la forme $K_*$, où $K_*$ désigne la famille des
\[
 K_*(\Delta^n)=K_{[n]}\qquad
 K_*\overset{\inserted{\text{par abus de notation}}}{=}(K_{[n]})_{n\geq0}
 \quad\inserted{\text{ou mieux }\Delta^n\longmapsto K_{[n]}}
\]
avec les opérations $1/2$ simpliciales entre elles. On fera attention qu'on écrit $K_{[n]}$ et non $K_n$, pour des raisons qui vont apparaître (\unsure{impérieuses|limpidement} lorsque $A$ est additive\ldots). De $=$ un objet de $A^*$ sera noté
\[
 K^*=(n\longmapsto K^{[n]})\overset{\text{abrév.}}{=}(K^{[n]})_n
\]

Nous aurons à travailler avec la situation où on a deux catégories $A$ et $B$, et une équivalence (notée \deleted{$\unsure{*}$} $M\longmapsto M^\vee$)
\[
 v:A^o\xrightarrow{\approx}B\qquad\text{d'où}\quad B^o\overset{\approx}{\unsure{\longrightarrow|-}}A
\]
[le plus souvent $A=B$ et $(\deleted{\unsure{M|\check{M}}}M^\vee)^\vee\underset{\substack{\text{isom}\\\text{fonct.}}}{\simeq}M$, avec compatibilités habituelles d'une autodualité\ldots]
on notera alors \inserted{souvent} par la $=$ lettre un objet de $A_*$ (ou $A^*$) et l'objet de $A^*$ (resp $A_*$) qui lui correspond par application de $v$, mais en indiquant la variance par la position du signe $*$
\[
 \begin{aligned}
 K^*&=(K_*)^\vee & K^{[n]}&=K_{[n]}^\vee\\
 K_*&=(K^*)^\vee & K_{[n]}&=K^{[n]\,\vee}
 \end{aligned}
\]

On s'intéressera surtout au cas où ($k$ étant un anneau fixé) \deleted{\unsure{$A=\text{la}$}}
\[
 \begin{aligned}
 \deleted{\text{cat. }A}\qquad A&=\deleted{\unsure{\mathrm{Ab}_k}}\ k\text{-modules à g. proj. de t.f.}\\
 B&=k\text{-modules à droite proj. de t.f.}
 \end{aligned}
\]
(si $k$ est commutatif, on a $A=B$ et on a une autodualité \unsure{sans|$v$ avec|dans|sur} $A$).

On a le foncteur canonique pl. fid
\[
 \mathbb{\Delta}\lhook\joinrel\longrightarrow\widehat{\mathbb{\Delta}}=\mathrm{Ss}\tag*{(*)}
\]
l'image par ce foncteur de $\Delta^n$ est noté $\Delta^n_*$ \inserted{ou mieux $\Delta^{[n]}_*$ (puisqu'il est covariant en $n$)}. On a donc
\[
 \Delta^n_*\ (\text{ou }\Delta^{[n]}_*)=
 \bigl(m\longmapsto\Delta^n_{[m]}=\mathrm{Hom}_{\mathrm{Ens}}(\Delta^m,\Delta^n)\bigr)
\]
et pour $\deleted{\unsure{F}}X_*\in\mathrm{Ob}\,\widehat{\mathbb{\Delta}}$,
\[
 \mathrm{Hom}(\Delta^n_*,\deleted{\unsure{F}}X_*)\simeq X_*(\Delta^n)\simeq X_{[n]}
\]
en particulier
\[
 \mathrm{Hom}(\Delta^n_*,\Delta^m_*)\simeq\Delta^m_*(\Delta^n_*)\simeq
 \underset{\substack{\Vert\\(\mathrm{Hom}_{\mathrm{ord}}(\Delta^n,\Delta^m))}}{\mathrm{Hom}_{\mathbb{\Delta}}(\Delta^n,\Delta^m)}
\]
Pour $n$ variable, $\Delta^n_*\in\widehat{\mathbb{\Delta}}$ dépend de façon covariante de $n$, \deleted{\unsure{i.e. o.|i.e.}} donc le foncteur $n\longmapsto\Delta^n_*$ (qui n'est autre que $(*)$) est un objet canonique de $\widehat{\mathbb{\Delta}}^*$, qu'on appelle \uline{objet cosimplicial universel} et note $\Delta^*$ (ou mieux $\Delta^*_* $)
\[
 \begin{aligned}
 \Delta^*_*\in\widehat{\mathbb{\Delta}}^* &=(\mathrm{Ens}_*)^*
 \quad\text{d'où }\forall n\geq0\\
 \Delta^{[n]}_*\in\widehat{\mathbb{\Delta}} \qquad
 \Delta^{[n]}_*&=\Delta^n_* =m\longmapsto
 \underset{\substack{\Vert\\\mathrm{Hom}_{\mathbb{\Delta}}(\Delta^m,\Delta^{\unsure{m|n}})}}{\Delta^n_{[m]}}
 \end{aligned}
\]

\deleted{\unsure{Ainsi}} [Une parenthèse de notations : Considérons une catégorie de la forme $(A_*)^*$, un objet de la dite est donc noté $K^*=\unsure{(|[}K^{[n]})_n$, où les $K^{[n]}$ sont des $\in A_*$, donc s'écrivent $K^{[n]}_*$,
\[
 K^{[n]}_*=(K^{[n]}_{[m]})_{m\geq0}
\]
Ainsi, on en vient à noter
\[
 K^*_*\quad\text{les objets de }(A_*)^*\simeq
 \underline{\mathrm{Hom}}(\mathbb{\Delta}\times\mathbb{\Delta}^o,A\unsure{)}
\]
\deleted{\unsure{on}}
% archive sheet number: 304
